\documentclass[11pt]{article}

\usepackage[a4paper,margin=28mm]{geometry}
\usepackage[T1]{fontenc}
\usepackage{lmodern}
\usepackage{microtype}
\usepackage{amsmath,amssymb,mathtools}
\usepackage{booktabs,longtable,array}
\usepackage{enumitem}
\usepackage{xcolor}
\usepackage{hyperref}

\hypersetup{
  colorlinks=true,
  linkcolor=blue!55!black,
  urlcolor=blue!55!black,
  pdftitle={Temporal Prolog: Normative Language Specification},
  pdfauthor={Temporal Prolog project}
}

\setlist{nosep}
\newcommand{\prev}{\mathbin{\bullet}}
\newcommand{\alwayspast}{\mathbin{\blacksquare}}
\newcommand{\once}{\mathbin{\blacklozenge}}
\newcommand{\eventuallypast}{\mathbin{\lozenge}}
\newcommand{\alwaysfuture}{\mathbin{\square}}
\newcommand{\sem}[1]{\llbracket #1 \rrbracket}
\newcommand{\FV}{\operatorname{fv}}
\newcommand{\true}{\mathsf{true}}
\newcommand{\false}{\mathsf{false}}
\newcommand{\at}{\mathsf{at}}
\newcommand{\since}{\mathbin{\mathsf{since}}}
\newcommand{\after}{\mathbin{\mathsf{after}}}
\newcommand{\until}{\mathbin{\mathsf{until}}}
\newcommand{\atnext}{\mathbin{\mathsf{atnext}}}

\title{Temporal Prolog\\\large Normative Language Specification and 1986 Paper Errata}
\author{Temporal Prolog project}
\date{Version 1.0 --- 21 July 2026}

\begin{document}
\maketitle

\begin{abstract}
This document specifies Temporal Prolog as implemented by this project.  Its
source is Takashi Sakuragawa's 1986 paper, but this specification fills in the
paper's intentionally informal concrete syntax, makes every temporal operator
denotationally explicit, and corrects inconsistencies identified in the
published transformation.  The specification separates the mathematical
language from the finite, range-restricted executable profile.  Both the
Haskell and Rust implementations are required to accept the same profile,
normalize it to the same semantics (up to fresh-name alpha equivalence), and
produce the same set of minimal worlds.
\end{abstract}

\tableofcontents

\section{Status and conformance}

The key words \textbf{MUST}, \textbf{MUST NOT}, \textbf{SHOULD}, and
\textbf{MAY} are normative.  A conforming implementation has two layers.

\begin{description}
  \item[Core semantics.] The trace and minimal-model semantics in
  Sections~\ref{sec:trace-semantics}--\ref{sec:model}.  This layer is the
  meaning of a program, including programs with recursion through negation.
  \item[Executable profile.] The decidable operational subset in
  Section~\ref{sec:profile}.  Implementations MUST reject inputs outside the
  profile rather than silently assign them a different meaning.
\end{description}

The words ``the paper'' below refer to Sakuragawa, \emph{Temporal Prolog},
RIMS K\={o}ky\={u}roku 586 (1986), pp.~305--329.

\section{Lexical and concrete syntax}

\subsection{Lexical classes}

A variable begins with an uppercase ASCII letter and continues with ASCII
letters, digits, or underscores.  A name begins with a lowercase ASCII letter
or underscore and has the same continuation characters.  Decimal integer
literals may have a leading minus sign and denote arbitrary-precision
mathematical integers.  Their spelling is canonicalized by numeric value, so
\verb|007| denotes the same term as \verb|7| and \verb|-0| denotes
\verb|0|.  A comment begins with \verb|%| and extends to the end of the line.
Whitespace and comments separate tokens.  Identifiers containing non-ASCII
characters MUST be rejected; Unicode is accepted only for the operator aliases
listed below.

Keyword reservation is contextual because predicates and constructors occupy
separate namespaces.  The formula-control words \verb|since|, \verb|after|,
\verb|for|, \verb|until|, \verb|atnext|, \verb|always|,
\verb|eventually|, and \verb|next| MUST NOT be used as predicate or pattern-
function names, but MAY be used as constructor names in term position.  Names
whose special meaning belongs to only one namespace are not globally reserved.
In particular, \verb|div| and \verb|mod| MAY be predicate names even though
their arity-two term forms are arithmetic, while external predicates such as
\verb|true| and \verb|is| MAY also occur as constructor names inside terms.
The call spellings \verb|true()| and \verb|false()| are equivalent to their
bare forms, and \verb|is(L,R)| is equivalent to \verb|L is R|.

The namespaces below are disjoint in the abstract language:
\[
V,\quad SF,\quad PF_i,\quad PF_e,\quad P_i,\quad P_e.
\]
They denote variables, constructor (Skolem) functions, internal and external
pattern functions, and internal and external predicates.  Each non-variable
symbol has a fixed arity.  Constants are arity-zero constructor functions.
The external predicates include $=$, \verb|is|/2, the four integer
comparisons, $\true$, $\false$, and $\at/1$.

The concrete implementation infers internal pattern-function names from
reduction heads.  Other callable names are predicates; constructor names are
identified by term position.  Predicate and constructor signatures are tracked
separately, so the same spelling MAY have one predicate arity and one
constructor arity.  Within either namespace, a program MUST use a symbol at one
fixed arity.

Every internal pattern-function name has one input arity $k$ across all of its
reductions.  Its term form has arity $k$, while its relational atom form has
arity $k+1$, with the final argument denoting the result.  Once a reduction
declares a pattern-function spelling, other uses of that spelling MUST obey
these two signatures.  External predicates have the fixed signatures listed
above and MAY occur only as conditions; they MUST NOT occur in rule results or
be inserted into a world as external input.  The arithmetic symbols $+$, $-$,
$*$, \verb|div|, and \verb|mod| have fixed arity two in term position.  An
implementation MUST validate these constraints before normalization and reject
the whole source program on a violation.  The same signatures remain fixed
across runtime queries, pending assertions, and every previously computed
world; runtime input MUST NOT introduce a conflicting arity.

The concrete spelling \verb|prefix_auxN| is not reserved: a source predicate
with that spelling remains a source predicate and MUST remain observable.
Implementations therefore track generated predicates by normalization
provenance rather than classifying names lexically.  When \texttt{N} is a
nonempty decimal suffix, fresh-name allocation interprets it as an
arbitrary-precision natural number and seeds the global counter strictly above
the greatest such source suffix.  This calculation MUST NOT truncate, ignore,
or overflow a suffix based on the host machine's integer or pointer width.

\subsection{Grammar}

The following EBNF is normative for the portable ASCII syntax.  Unicode
aliases are listed in Table~\ref{tab:aliases}.

\begin{verbatim}
program      ::= item*
item         ::= rule "." | reduction "."
rule         ::= result | condition "=>" result
reduction    ::= call "->" term

result       ::= result-and
                 [("until" | "atnext") condition]
result-and   ::= result-unary ("/\\" result-unary)*
result-unary ::= "always" result-unary
               | "next" result-unary
               | result-atom
result-atom  ::= atom | reduction | "(" result ")"

condition       ::= condition-and
                    [("since" | "after") condition-and
                     | "for" positive-integer]
condition-and   ::= condition-unary ("/\\" condition-unary)*
condition-unary ::= ("~" | "@" | "#" | "?") condition-unary
                  | "eventually" condition-unary
                  | condition-atom
condition-atom  ::= atom | "(" condition ")"

atom         ::= call | name
               | term ("=" | "!=" | "\\=" | "is"
                       | ">" | "<" | ">=" | "<=") term
call         ::= name "(" [term ("," term)*] ")"

term         ::= variable | integer | name | call
               | "[" [term ("," term)* ["|" term]] "]"
               | "@" term
               | term ("+" | "-" | "*" | "div" | "mod") term
               | "(" term ")"
\end{verbatim}

The \texttt{positive-integer} metavariable uses the same arbitrary-precision
decimal domain as an integer term, but excludes zero and negative values.

Unary operators bind most tightly.  Condition conjunction binds more tightly
than \verb|since|, \verb|after|, and \verb|for|.  Those three binary past-time
operators are non-associative; chained uses require parentheses.  Within a result,
\verb|always| and \verb|next| bind most tightly, result conjunction binds next,
and \verb|until| and \verb|atnext| bind least tightly.  The latter two operators
are non-associative: nested uses require parentheses.  Thus
\verb|p /\ q until stop| means \verb|(p /\ q) until stop|.  Implication binds
least tightly overall.

\begin{table}[ht]
\centering
\begin{tabular}{lll}
\toprule
ASCII & Unicode & Meaning \\
\midrule
\verb|=>| & $\Rightarrow$ & implication \\
\verb|/\| & $\wedge$ & conjunction \\
\verb|->| & $\to$ & reduction \\
\verb|@| & $\prev$ & previous time \\
\verb|~| & $\neg$ & negation \\
\verb|#| & $\alwayspast$ & continuously in the past \\
\verb|?| & $\once$ & once in the past \\
\verb|eventually| & $\eventuallypast$ & once in the past \\
\verb|always| & $\alwaysfuture$ & henceforth \\
\verb|next| & $\bigcirc$ & next time \\
\bottomrule
\end{tabular}
\caption{Portable ASCII spellings and accepted Unicode aliases.}
\label{tab:aliases}
\end{table}

Both \texttt{U+25CF} BLACK CIRCLE and \texttt{U+2022} BULLET are accepted
aliases for \verb|@|.  \texttt{U+25C6} BLACK DIAMOND is the Unicode spelling of
\verb|?|, while \texttt{U+25C7} WHITE DIAMOND is the Unicode spelling of
\verb|eventually|.  Implementations MUST preserve those operator families in
the parsed source AST, even though \verb|?| and \verb|eventually| normalize to
the same past-existential semantics.

\subsection{Well-formed terms and formulas}

Terms are generated by
\[
t ::= X \mid f(t_1,\ldots,t_k) \mid \prev t.
\]
The paper restricts $\prev t$ to a term containing at least one occurrence of
a pattern function.  The executable profile enforces that restriction.  In
particular, \verb|@X| is not a portable way to refer to a variable's old
binding.  Previous-time pattern-function values are defined in
Section~\ref{sec:pf}.

An atom is $p(t_1,\ldots,t_k)$.  A condition is an atom closed under negation,
previous time, past operators, and conjunction.  A result is an internal atom
or internal reduction closed under implication, conjunction, and future
operators.  Free variables in each top-level item are universally quantified.

\section{Traces and temporal truth}\label{sec:trace-semantics}

A trace is an infinite sequence $M=\langle w_0,w_1,\ldots\rangle$.  Each
$w_n$ is a set of ground atoms.  Constructor interpretations and the Herbrand
domain are time invariant; predicate truth may vary with time.  We write
$M,n,\rho\models c$ for condition $c$ at time $n$ under grounding
environment $\rho$.

Atomic truth, conjunction, and negation are classical at a fixed world:
\begin{align*}
M,n,\rho &\models a &&\Longleftrightarrow \rho(a)\in w_n,\\
M,n,\rho &\models \neg c &&\Longleftrightarrow M,n,\rho\not\models c,\\
M,n,\rho &\models c\wedge d &&\Longleftrightarrow
  M,n,\rho\models c\text{ and }M,n,\rho\models d.
\end{align*}

Previous time is false before the trace begins, irrespective of its operand:
\[
M,n,\rho\models \prev c
\Longleftrightarrow n>0\text{ and }M,n-1,\rho\models c.
\]
Consequently $\prev\neg p$ is false at $w_0$, while
$\neg\prev p$ is true there.  An implementation MUST NOT commute negation
through $\prev$ at the initial world.

The derived past operators have the following meanings:
\begin{align*}
M,n,\rho\models \alwayspast c
&\Longleftrightarrow \forall i\ (0\le i\le n\Rightarrow M,i,\rho\models c),\\
M,n,\rho\models \once c
&\Longleftrightarrow \exists i\ (0\le i\le n\land M,i,\rho\models c),\\
M,n,\rho\models c\since d
&\Longleftrightarrow \exists j\le n\,.
  M,j,\rho\models d\land
  \forall k\,(j\le k\le n\Rightarrow M,k,\rho\models c),\\
M,n,\rho\models c\after d
&\Longleftrightarrow \exists i,j\,.
  0\le i<j\le n\land M,i,\rho\models d\land M,j,\rho\models c,\\
M,n,\rho\models c\ \mathsf{for}\ k
&\Longleftrightarrow k>0\land n\ge k-1\land
  \forall i<k\,. M,n-i,\rho\models c.
\end{align*}
The \verb|after| relation is strict and, once witnessed, remains true.  This
is the paper's prose definition; it corrects the inconsistent printed
Step~2(4), as explained in Section~\ref{sec:errata}.

The built-in $\at(k)$ is true exactly at $w_k$.  $\true$ is always true and
$\false$ always false.  Equality is syntactic equality in the Herbrand
universe.

\section{Result obligations}

A top-level item is required to hold at every time.  We define
$M,n,\rho\models_R r$ as follows.

An atomic result $a$ requires $M,n,\rho\models a$.  A reduction result is
defined relationally in Section~\ref{sec:pf}.  Conjunction requires both
results.  Implication has the usual meaning:
\[
M,n,\rho\models_R c\Rightarrow r
\Longleftrightarrow M,n,\rho\not\models c\text{ or }M,n,\rho\models_R r.
\]

Future result operators create obligations from the evaluation point:
\begin{align*}
M,n,\rho\models_R \alwaysfuture r
&\Longleftrightarrow \forall k\ge n\,. M,k,\rho\models_R r,\\
M,n,\rho\models_R \bigcirc r
&\Longleftrightarrow M,n+1,\rho\models_R r.
\end{align*}
For $r\until c$, let $j$ be the least $k\ge n$ satisfying $c$, if one
exists.  The result $r$ is required at every $k$ with $n\le k<j$; if no such
$j$ exists, $r$ is required forever from $n$.  For $r\atnext c$, $r$ is
required at that least $j$ and no obligation is imposed before it.  If $c$
never occurs, \verb|atnext| imposes no result obligation.

A program $P$ is satisfied by $M$ iff every universally grounded top-level
item holds at every $n\ge0$.

\section{Normalization}\label{sec:normalization}

Normalization produces clauses
\[
\prev^{m_1}(\neg)c_1\wedge\cdots\wedge
\prev^{m_k}(\neg)c_k\Rightarrow d,
\]
where each $c_i$ is atomic, each $m_i\ge0$, and $d$ is an internal atom.
Fresh predicates and variables MUST be absent from the source program.  Rules
below are equivalences under the minimal-model semantics; fresh names are
local to a normalization run.

\subsection{Step 1: future results}

Result conjunction distributes into separate clauses.  Bare future formulas
simplify as
\[
\alwaysfuture q\rightsquigarrow q,
\qquad
q\until a\rightsquigarrow \neg a\Rightarrow q,
\qquad
q\atnext a\rightsquigarrow a\Rightarrow q.
\]

For a conditional always result, with fresh $p(\vec X)$ and
$\vec X=\FV(q)$:
\begin{align*}
a\Rightarrow\alwaysfuture q\quad\rightsquigarrow\quad
&a\Rightarrow p(\vec X),\\
&\prev p(\vec X)\Rightarrow p(\vec X),\\
&p(\vec X)\Rightarrow q.
\end{align*}

For conditional until, let $\vec X=\FV(q,b)$:
\begin{align*}
a\Rightarrow q\until b\quad\rightsquigarrow\quad
&a\Rightarrow p(\vec X),\\
&\prev p(\vec X)\wedge\prev\neg b\Rightarrow p(\vec X),\\
&p(\vec X)\wedge\neg b\Rightarrow q.
\end{align*}

Conditional \verb|atnext| uses the same first two clauses and replaces the
third with $p(\vec X)\wedge b\Rightarrow q$.  The project extension
$a\Rightarrow\bigcirc q$ becomes
$a\Rightarrow p(\FV(q))$ and $\prev p(\FV(q))\Rightarrow q$.

\subsection{Step 2: past conditions}

For each occurrence below, replace the occurrence by $p$ in its surrounding
clause.  Predicate arguments are exactly the displayed free-variable set,
not the variables of the surrounding clause.

\begin{center}
\begin{longtable}{>{\raggedright\arraybackslash}p{0.18\linewidth}
                  >{\raggedright\arraybackslash}p{0.70\linewidth}}
\toprule
Occurrence & Defining clauses \\
\midrule
$\alwayspast a$ &
$a\wedge\at(0)\Rightarrow p(\FV(a))$;
$\prev p(\FV(a))\wedge a\Rightarrow p(\FV(a))$ \\
$\once a$ &
$a\Rightarrow p(\FV(a))$;
$\prev p(\FV(a))\Rightarrow p(\FV(a))$ \\
$a\since b$ &
$b\wedge a\Rightarrow p(\FV(a,b))$;
$\prev p(\FV(a,b))\wedge a\Rightarrow p(\FV(a,b))$ \\
$a\after b$ &
$b\Rightarrow s(\FV(b))$;
$\prev s(\FV(b))\Rightarrow s(\FV(b))$;
$\prev s(\FV(b))\wedge a\Rightarrow p(\FV(a,b))$;
$\prev p(\FV(a,b))\Rightarrow p(\FV(a,b))$ \\
$a\ \mathsf{for}\ k$ &
replace with $a\wedge\prev a\wedge\cdots\wedge\prev^{k-1}a$, where $k>0$ \\
\bottomrule
\end{longtable}
\end{center}

The auxiliary $s$ in \verb|after| records that the right operand occurred;
$p$ records a later left-operand witness.  Using $\prev s$ makes the order
strict.  Both predicates retain exactly the variables necessary to relate
their witnesses.

\subsection{Step 3: pattern functions}\label{sec:pf}

An internal reduction
\[
c\Rightarrow f(t_1,\ldots,t_k)\to t_0
\]
becomes the relational clause
\[
c\Rightarrow f(t_1,\ldots,t_k,t_0).
\]
The new predicate has arity $k+1$.  Multiple reductions denote a relation;
the language does not silently impose functional uniqueness.

For the outermost, rightmost pattern-function occurrence
$f(t_1,\ldots,t_k)$ inside a predicate term, introduce fresh $X$, replace the
occurrence by $X$, and conjoin $f(t_1,\ldots,t_k,X)$.  If the occurrence lies
under $m$ term-level previous operators, conjoin
$\prev^m f(t_1,\ldots,t_k,X)$ while leaving the enclosing predicate at its
original condition depth.  Repeat until no pattern function occurs in a
term, then erase every remaining term-level $\prev$ wrapper.

For example,
\[
\mathsf{present}(\prev\mathsf{lookup}(k))
\rightsquigarrow
\mathsf{present}(X)\wedge\prev\mathsf{lookup}(k,X).
\]
Moving \verb|present| into the previous world would be non-conformant.

\subsection{Steps 4 and 5: negation and previous time}

If $\neg a$ occurs and $a$ is not atomic, introduce
$a\Rightarrow p(\FV(a))$ and replace the occurrence by
$\neg p(\FV(a))$.  This rule is essential for $\neg\prev a$: directly
rewriting it as $\prev\neg a$ is invalid at $w_0$.

Finally distribute $\prev$ over conjunction.  A normal condition has exactly
the shape $\prev^m a$ or $\prev^m\neg a$.

\section{Models and nondeterminism}\label{sec:model}

Past worlds and external atoms at the current world are fixed inputs.  For a
normal program, remove all previous-time body atoms when constructing the
current dependency graph.  There is an edge $p\to q$ when a current-world
atom with predicate $q$ occurs in the body of a clause headed by predicate
$p$; retain whether the occurrence is positive or negative.

Collapse strongly connected components (SCCs) and order them so every
dependency SCC precedes its consumers.  Within a current world, compare two
models lexicographically by the sets of atoms in this SCC order: at the first
component where they differ, a model is smaller only when its component set
is a proper subset of the other's.  Incomparable component sets remain
incomparable.  A model is minimal when no model is smaller by this order.

If no SCC contains a negative edge, the program satisfies the paper's
condition~1 and has a unique least model.  It is computed SCC-by-SCC by least
fixed points.  Previous-time conditions do not create current-world
dependencies.

If a negative edge lies inside an SCC, a conforming implementation MUST NOT
pretend that the program is stratified.  Its semantics is the set of minimal
models.  An execution MAY choose any member, and an API intended for testing
SHOULD expose all members.  For example, the clauses printed in paper
Section~4.7 (made range restricted here by an explicit domain condition) are:
\begin{verbatim}
assign(X) /\ ~assigned_to_another(X) => assigned_to(X).
assigned_to(X) /\ assign(Y) /\ X != Y => assigned_to_another(Y).
\end{verbatim}
With simultaneous requests from 1 and 2, the paper claims two minimal worlds,
one assigning each requester.  Under its stated model order there is also a
third: the ``assigned to another'' blocker is true for both requesters and
neither requester is assigned.
Neither blocker can be removed without forcing an incomparable assignment,
so this world is minimal.  Implementations MUST expose all three for the
published program.  The following corrected encoding has exactly the intended
two minimal worlds:
\begin{verbatim}
assign(X) /\ assign(Y) /\ X != Y /\ ~assigned_to(Y)
  => assigned_to(X).
\end{verbatim}

The mathematical language calls a program legal when it has at least one
model for every interpretation of its external predicates.  This property is
not decidable for the unrestricted Herbrand universe.

\section{Executable profile}\label{sec:profile}

The Haskell and Rust implementations execute the following common profile.

\begin{enumerate}
  \item Assertions supplied for a step are ground and form part of the fixed
  external input for that step.  They MUST preserve predicate and constructor
  signatures established by the normalized program and all earlier or pending
  assertions.  A newly introduced input symbol establishes its signature for
  the remainder of the run.
  \item Assertions MUST NOT use an external-predicate name, an internal
  pattern-function relation, or a predicate introduced by normalization.
  Runtime atoms MUST NOT contain a term-level \verb|@| wrapper or a malformed
  arithmetic operator.
  \item Query atoms obey the same signature rules and MUST NOT address
  normalization-generated predicates.  Internal pattern-function relations
  MAY be queried at their relational arity.
  \item Every variable used by a negated condition or a forward-chained head
  is grounded by positive body conditions, $\at(X)$, unification with a
  ground term, or a ground arithmetic evaluation before it is observed.
  \item A parsed \verb|for| count is retained as an arbitrary-precision
  positive integer.  The portable executable profile admits counts through
  1,000 inclusive.  A larger count MUST be rejected as a normalization
  resource error before expansion; it MUST NOT wrap, truncate, or depend on
  the host machine's integer or pointer width.
  \item Steps 1, 2, and 4 each admit 1,000 productive rounds.  Reaching normal
  form on round 1,000 succeeds; a target afterward is a resource error.
  Implementations MUST NOT spend a round on a no-change convergence probe.
  \item Candidate generation for a world reaches a finite fixed point within
  the configured fuel limit.  A program that generates an unbounded Herbrand
  base is rejected for that run.
  \item The general evaluator's candidate base is closed under relaxed
  positive derivation and under every grounded current-world negative atom
  whose predicate is internal and whose remaining body is satisfied in that
  relaxed closure.  These steps repeat to a fixed point.  This second clause
  is necessary because the paper uses classical minimal models: an
  unsupported atom may occur solely to block a negative antecedent.
  \item General negative SCCs are executed only when the finite candidate
  base is no larger than the configured enumeration limit.  Exceeding the
  limit is a resource error, not a semantic result.
  \item Internal pattern-function relations are solved by alpha-renamed
  backward chaining.  Reaching its recursion limit is reported as a resource
  error in conformance mode; it MUST NOT be reported as logical failure.
\end{enumerate}

For condition-1 programs, the optimized stratified evaluator and the general
minimal-model evaluator MUST agree.  Differential tests compare both engines
and compare the Haskell and Rust implementations on canonicalized worlds.

\section{External operations and project extensions}

The portable executable profile defines arbitrary-precision integer arithmetic
terms $+$, $-$, $*$, \verb|div|, and \verb|mod|.  Multiplication,
\verb|div|, and \verb|mod| share one left-associative precedence level above
addition and subtraction.  The prefix forms \verb|div(a,b)| and
\verb|mod(a,b)| are equivalent functional spellings for the corresponding
arithmetic expressions.

For integers $a$ and nonzero $b$, $a\mathbin{\mathrm{div}}b$ is the quotient
$q=\lfloor a/b\rfloor$, and $a\mathbin{\mathrm{mod}}b$ is the unique remainder
$r=a-bq$.  Thus $r$ is zero or has the sign of $b$, and $|r|<|b|$.
The predicate $X\ \mathsf{is}\ E$ evaluates a ground expression $E$ and
unifies $X$ with its integer result.  Comparisons require ground integer
operands.  Division by zero and nonground or noninteger arithmetic fail
without binding and MUST NOT fall through to stored-predicate lookup.
\verb|!=| and \verb|\=| abbreviate negated Herbrand equality.
The remaining external predicates are Herbrand equality $=$, the four integer
comparisons, \verb|at|/1 for the current world number, and the constants
$\true$ and $\false$.  These operations are determined by the evaluator rather
than world membership, so an implementation MUST reject them in scheduled
assertions and model-checking input choices.  A public query API MUST evaluate
these predicates by the same rules used for program conditions and return
their substitutions; it MUST NOT treat a well-formed successful external query
as an absent stored fact.

The project accepts \verb|eventually c| as a spelling of $\once c$ and
\verb|next r| as the next-result operator.  These are conservative extensions
and are tested by both implementations.

\section{Published-paper errata and clarifications}\label{sec:errata}

\begin{longtable}{>{\raggedright\arraybackslash}p{0.17\linewidth}
                  >{\raggedright\arraybackslash}p{0.35\linewidth}
                  >{\raggedright\arraybackslash}p{0.38\linewidth}}
\toprule
Location & Published text or ambiguity & Normative resolution \\
\midrule
\endfirsthead
\toprule
Location & Published text or ambiguity & Normative resolution \\
\midrule
\endhead
Section 3, \verb|after|; Step 2(4) &
The prose requires a left-operand occurrence after a right-operand occurrence,
but the printed recurrence merely remembers the more recent operand and is
true after the left operand even if the right never occurred. &
Use the strict two-latch construction in Section~\ref{sec:normalization}.  A
same-world pair is not ``after,'' and a witnessed pair remains witnessed. \\

Section 3, \verb|for| &
The metavariable is first called non-negative, then explicitly positive. &
The operand MUST be positive.  Zero is rejected. \\

Step 1(2) and Step 2 auxiliary arguments &
Several displayed subscripts alternate $k$ and $l$, and an implementation
might capture variables from the surrounding clause. &
Use exactly the free-variable sets printed by the defining operator:
$\FV(q)$, $\FV(q,b)$, $\FV(a)$, or $\FV(a,b)$. \\

Step 3 term-level previous &
The paper says to erase term-level bullets after expansion but does not warn
that lifting the enclosing predicate changes its time. &
Only the generated pattern-function condition receives the occurrence's term
depth; the enclosing atom keeps its condition depth. \\

Steps 4--5 at $w_0$ &
The normal form places negation inside previous operators, which can tempt an
implementation to rewrite $\neg\prev a$ as $\prev\neg a$. &
Step 4 MUST introduce an auxiliary before distribution.  The two formulas
differ at $w_0$. \\

Section 5.3 &
Condition~1 is a sufficient least-model condition, not the definition of all
legal programs; Section 4.7 intentionally uses a negative SCC. &
Expose minimal-model nondeterminism for finite executable cases.  Rejecting
every negative SCC is a supported fast-path limitation only when clearly
reported, not a conformant semantic answer. \\

Section 4.7 &
The nondeterministic assignment program is claimed to have two choices, but
its two unsupported \texttt{assigned\_to\_another} atoms form a third
minimal model in which nobody is assigned. &
Retain all three models for the published clauses, as required by the formal
minimal-model semantics.  Use the corrected direct-choice clause in
Section~\ref{sec:model} when exactly two assignment outcomes are intended. \\

Concrete syntax and external domains &
The paper intentionally gives no machine grammar or domain-declaration syntax. &
Use Section 2's grammar and Section~\ref{sec:profile}'s ground assertion and
range-restriction rules. \\
\bottomrule
\end{longtable}

\section{Conformance requirements}

A release is conformant only when all of the following are true.

\begin{enumerate}
  \item Both implementations parse every shared positive corpus case and
  reject every shared negative case.  Rejections cover symbol-arity conflicts,
  non-ASCII identifiers, malformed pattern-function and external-operation
  signatures, oversized \verb|for| counts, and external predicates used as
  results or assertions.  Positive cases
  exercise contextual keyword reservation across predicate and constructor
  namespaces.  Runtime rejection
  cases cover signature changes across source,
  pending input, and history; internal-name injection; term-level previous;
  and malformed arithmetic.  Public-query tests cover every external predicate
  family and its variable bindings.
  \item Normalization invariants hold: no surface temporal operator or
  term-level previous wrapper remains, every body condition is normal, and
  fresh identifiers do not collide with source identifiers.  Generated-name
  provenance is exact: source identifiers resembling the fresh-name spelling
  remain observable, while only actual auxiliaries are hidden by default.
  Arbitrary-precision source suffixes seed the same exact next fresh name in
  both implementations.
  \item Canonical sets of worlds agree for every shared scenario, including
  initial-world negation, corrected \verb|after|, recursive pattern
  functions, arbitrary-precision signed arithmetic, unsupported classical
  blockers, and Section 4.7 nondeterminism.
  \item The stratified and general evaluators agree wherever condition~1
  holds.
  \item Benchmarks report workload parameters, iterations, elapsed time, and
  a semantic digest.  A timing result without a checked digest is invalid.
\end{enumerate}

\begin{thebibliography}{9}
\bibitem{sak1986}
Takashi Sakuragawa.
\newblock Temporal Prolog.
\newblock \emph{RIMS K\={o}ky\={u}roku}, 586:305--329, 1986.

\bibitem{clark1977}
Keith L. Clark.
\newblock Negation as failure.
\newblock In \emph{Logic and Data Bases}, pages 293--322, 1977.
\end{thebibliography}

\end{document}
